% Part: sets-functions-relations
% Chapter: arithmetization
% Section: checking-details
%
\documentclass[../../../include/open-logic-section]{subfiles}


\begin{document}
	
\olfileid{sfr}{arith}{check}
\olsection{有序环与域}

在本章中，我们一直声称某些定义“表现得理应如此”。在这个技术附录里，我们将阐明其确切含义，并（概述如何）证明这些定义确实表现得“正确”。

在 \olref[int]{sec} 中，我们定义了 $\Int$ 上的加法和乘法。我们希望证明，按我们的定义，它们赋予了 $\Int$ 我们“希望”它具有的结构。具体来说，这个结构就是交换环的结构。


\begin{defn}
	一个\emph{交换环}是一个集合 $S$，配备有特定元素 $0$ 和 $1$ 以及运算 $+$ 和 $\times$，满足以下八个公式：
	\begin{align*}
		\emph{Associativity}&&a + (b+ c) & = (a + b) + c \\
		&& (a \times b) \times c & = a \times (b\times c)\\
		\emph{Commutativity}&&a + b &= b+ a  \\
		&&  a \times b&= b\times a\\ 
		\emph{Identities}&&a + 0 &= a \\
		&& a \times 1 &= a\\
		\emph{Additive Inverse}&&(\exists b\in S)0&=a + b\\
		\emph{Distributivity}&&a \times (b+ c ) &= (a \times b) + (a \times c)
	\end{align*}
	隐含地，这些公式的全称量词辖域均限制于 $S$。并且注意，这里的元素 $0$ 和~$1$ 未必是同名的自然数。
\end{defn}


所以，要验证整数构成一个交换环，我们只需要检查它满足这八个条件。这些条件都不难验证，但有些繁琐。例如，以下是加法\emph{结合律}的证明：

\begin{proof}
取定 $i, j, k \in \Int$。于是存在 $a_1, b_1, a_2, b_2, a_3, b_3 \in
\Nat$ 使得 $i = \equivrep{a_1, b_1}{}$、$j =
\equivrep{a_2,b_2}{}$ 与 $k = \equivrep{a_3, b_3}{}$ 成立。
（为了可读性，我们写作“$\equivrep{x, y}{}$”而非“$\equivrep{x, y}{\Intequiv}$”；本节中我们都会这样做。）现在：
\begin{align*}
	i + (j + k) &= \equivrep{a_1, b_1}{}+(\equivrep{a_2, b_2}{} + \equivrep{a_3, b_3}{}) \\
	&= \equivrep{a_1,  b_1}{} + \equivrep{a_2+a_3, b_2+b_3}{}\\
	&= \equivrep{a_1 + (a_2 + a_3), b_1 + (b_2 + b_3)}{}\\
	&= \equivrep{(a_1 + a_2) + a_3, (b_1 + b_2) + b_3}{}\\
	&= \equivrep{a_1 + a_2, b_1 + b_2}{} + \equivrep{a_3, b_3}{}\\
	&= (\equivrep{a_1, b_1}{} + \equivrep{a_2, b_2}{}) + \equivrep{a_3, b_3}{}\\
	&= (i+j) + k
\end{align*}
这里我们自由运用了 $\Nat$ 上加法的性质。
\end{proof}



同样地，以下是\emph{加法逆元}的证明：

\begin{proof}
取定 $i \in \Int$，即存在 $a,b \in \Nat$ 使得 $i = \equivrep{a,b}{}$。令 $j = \equivrep{b,a}{} \in \Int$。利用自然数的已知性质，$(a+b) + 0 = 0 + (a+b)$，于是根据定义有 $\tuple{a+b, b+a} \sim_\Int \tuple{0,0}$，因此 $\equivrep{a+b, b+a}{} = \equivrep{0, 0}{} = 0_\Int$。这样我们就得到 $i + j =
\equivrep{a,b}{}+\equivrep{b,a}{}=\equivrep{a+b, b+a}{}= \equivrep{0,
0}{} = 0_\Int$。
\end{proof}

以下是\emph{分配律}的一个证明：
\begin{proof}
如上所述，取定 $i = \equivrep{a_1, b_1}{}$、$j =
\equivrep{a_2,b_2}{}$ 和 $k = \equivrep{a_3, b_3}{}$。现在：
\begin{align*}
	i \times (j + k) 
	&= \equivrep{a_1, b_1}{} \times (\equivrep{a_2,b_2}{} + \equivrep{a_3, b_3}{})\\
	&= \equivrep{a_1, b_1}{} \times \equivrep{a_2 + a_3,b_2+b_3}{}\\
	&= \equivrep{a_1  (a_2 + a_3) + b_1  (b_2+b_3), a_1  (b_2 + b_3) + b_1 (a_2 + a_3)}{}\\
	&= \equivrep{a_1 a_2 + a_1a_3 + b_1 b_2+b_1b_3, a_1 b_2 + a_1b_3 + a_2b_1 + a_3b_1}{}\\		
%		&= \equivrep{a_1 a_2 + b_1b_2 + a_1a_3 + b_1 b_2+b_1b_3, a_1 b_2 + a_1b_3 + a_2b_1 + a_3b_1}{}\\		
	&= \equivrep{a_1a_2 + b_1b_2, a_1b_2 + a_2b_1}{} + \equivrep{a_1a_3 + b_1b_3, a_1b_3 + a_3b_1}{}\\
	&= (\equivrep{a_1, b_1}{} \times \equivrep{a_2,b_2}{}) + (\equivrep{a_1, b_1}{} \times  \equivrep{a_3, b_3}{})\\
	&= (i \times j) + (i \times k)
\end{align*}
\end{proof}

我们把证明其余五个条件留作练习。
完成之后，我们就证明了 $\Int$ 构成一个交换环，也就是说，（按上述定义的）加法和乘法行为符合预期。

\begin{prob}
证明 $\Int$ 是一个交换环。
\end{prob}

但我们的任务还没完成。除了定义 $\Int$ 上的加法和乘法，我们还定义了一个序关系 $\leq$，我们必须验证其行为是否符合预期。更具体地说，我们必须证明 $\Int$ 构成一个\emph{有序}环。\footnote{回顾 \olref[sfr][rel][ord]{def:linearorder}，全序是一个自反、传递、反对称且连通的关系。在序关系的语境下，连通性有时被称为\emph{三分律}，因为对任意 $a$ 和 $b$，我们有 $a \leq b \lor a
	= b \lor a \geq b$。}

\begin{defn}
一个\emph{有序环}是一个交换环，并且它还配备了一个全序关系 $\leq$，满足：
\begin{align*}
	a \leq b &\lif a + c \leq b + c\\
	(a \leq b \land 0 \leq c) &\lif a \times c \leq b \times c
\end{align*}
\end{defn}

\begin{prob}
证明 $\Int$ 是一个有序环。
\end{prob}

和之前一样，证明我们构造的~$\Int$ 是一个有序环，这个过程是繁琐但常规的。我们将把这个问题留给你。

这样整数就处理完了。但现在我们需要对有理数证明非常类似的结论。具体来说，我们现在需要证明，在我们给定的 $+$、$\times$ 和 $\leq$ 的定义下，有理数构成一个有序\emph{域}：

\begin{defn}\ollabel{orderedfield}
一个\emph{有序域}是一个同时满足以下条件的有序环：
\begin{align*}
	\emph{Multiplicative Inverse}& & (\forall a \in S \setminus \{0\})(\exists b \in S) a\times b& = 1
\end{align*}
\end{defn}

一旦你证明了 $\Int$ 构成一个有序环，那么证明 $\Rat$ 构成一个有序域就容易（但依旧繁琐）了。

\begin{prob}
证明 $\Rat$ 是一个有序域。
\end{prob}

处理完整数和有理数，就只剩下处理实数了。具体来说，我们需要证明 $\Real$ 构成一个\emph{完备的}有序域，即具有完备性质的有序域。现在，\olref[cuts]{realcompleteness} 已经证明了 $\Real$ 具有完备性质。然而，我们仍需进行验证 $\Real$ 是一个有序域的（繁琐）工作。

在进行\emph{那项}繁琐的练习之前，我们需要检查一些更“直接”的事情。例如，我们需要保证，对于任何分割 $\alpha$ 和 $\beta$，按定义给出的 $\alpha + \beta$ 确实是一个\emph{分割}。以下是该事实的证明：

\begin{proof}
由于 $\alpha$ 和 $\beta$ 都是分割，$\alpha + \beta = \Setabs{p
+ q}{p \in \alpha \land q \in \beta}$ 是 $\Rat$ 的一个非空真子集。
现在，设 $x < p + q$，其中 $p \in \alpha$ 且 $q \in \beta$。
那么 $x - p < q$，所以 $x - p \in \beta$，并且 $x = p + (x - p) \in
\alpha + \beta$。
因此 $\alpha + \beta$ 是 $\Rat$ 的一个初始段。
最后，对任意 $p + q \in \alpha + \beta$，由于 $\alpha$ 和
$\beta$ 都是分割，存在 $p_1 \in \alpha$ 和 $q_1 \in \beta$
使得 $p < p_1$ 且 $q < q_1$；于是 $p + q < p_1 + q_1 \in \alpha +
\beta$；所以 $\alpha + \beta$ 没有最大元。
\end{proof}

类似的工作将使你能够验证 $\alpha - \beta$、$\alpha \times \beta$ 和 $\alpha \div \beta$ 也是分割（对于最后一种情况，需忽略 $\beta$ 是零分割的情形）。不过，我们还是直接把这些留给你。

\begin{prob}
证明 $\Real$ 是一个有序域。
\end{prob}

但这里还有一个小的遗留问题需要处理。在
\olref[cuts]{sec}中，我们建议可以取 $\sqrt{2} =
\Setabs{p \in \Rat}{p < 0 \text{ or }p^2 < 2}$。但我们确实需要证明
这个集合是一个\emph{分割}。以下就是这个事实的证明：

\begin{proof}
显然这是有理数的一个非空真初始段；因此
只需证明它没有最大元。具体来说，只需证明：
若 $p$ 是满足 $p^2 < 2$ 的正有理数且 $q =
\frac{2p+2}{p+2}$，则有 $p < q$ 且 $q^2 < 2$。
要说明 $p < q$，只需注意：
\begin{align*}
	p^2 &< 2\\
	p^2 + 2p &< 2 + 2p\\
	p(p + 2) &< 2 + 2p\\
	p &< \tfrac{2+2p}{p+2} = q
\end{align*}
要说明 $q^2 < 2$，只需注意：
\begin{align*}
	p^2 &< 2\\
	2p^2 + 4p + 2 &< p^2 + 4p+ 4\\
	4p^2 + 8p + 4 &< 2(p^2 + 4p + 4)\\
	(2p+2)^2 & <2(p+2)^2\\
	\tfrac{(2p+2)^2}{(p+2)^2} &< 2\\
	q^2 &< 2
\end{align*}
\end{proof}

\end{document}